% Applied Sciences MDPI template
\documentclass[journal,article,submit,pdftex,moreauthors]{Definitions/mdpi}

% MDPI internal commands
\firstpage{1}
\makeatletter
\setcounter{page}{\@firstpage}
\makeatother
\pubvolume{1}
\issuenum{1}
\articlenumber{0}
\pubyear{2026}
\copyrightyear{2026}
\datereceived{ }
\daterevised{ }
\dateaccepted{ }
\datepublished{ }

% Title
\Title{Design and Evaluation of an Open-Source LLM-Based Intelligent Tutoring System for Programming Education}

% Authors
\Author{Firstname Lastname $^{1}$* and Firstname Lastname $^{2}$}

\AuthorNames{Firstname Lastname and Firstname Lastname}

\address{%
$^{1}$ \quad Department of Computer Science, University Name, City, Country; e-mail@e-mail.com\\
$^{2}$ \quad Department of Computer Science, University Name, City, Country; e-mail@e-mail.com}

\corres{Correspondence: e-mail@e-mail.com}

% Abstract
\abstract{The integration of Large Language Models (LLMs) into educational technology has opened new possibilities for personalized programming instruction. This paper presents CodeTutor-ITS, an intelligent tutoring system powered by the open-source Qwen2.5-7B-Instruct model for programming education. The system features four core modules: (1) a dialogue-based tutoring engine that provides guided instruction through Socratic questioning, (2) a multi-level adaptive hint system that progressively reveals information based on student struggles, (3) an exercise generator that creates contextually appropriate programming problems, and (4) a knowledge tracking module that monitors student mastery across programming concepts. Unlike proprietary LLM-based solutions, our system runs entirely locally, ensuring data privacy and zero operational cost. We evaluate the system through four experiments: tutoring quality assessment using BLEU and ROUGE-L metrics on 50 programming questions, a prompt strategy comparison (zero-shot vs. few-shot vs. chain-of-thought), an ablation study measuring the contribution of each module, and an LLM-as-Judge evaluation using Qwen2.5-72B as an automated evaluator across four quality dimensions. Results demonstrate that the chain-of-thought prompting strategy achieves the highest completeness score of 0.80, and the ablation study shows that removing both knowledge tracking and adaptive hints reduces success rate from 100\% to 58\%. The system achieves an overall LLM-as-Judge score of 4.51 out of 5.0, with particularly strong performance in accuracy (4.70) and educational value (4.64). The system provides a practical, cost-effective solution for programming education institutions seeking to leverage AI-powered tutoring without dependence on commercial APIs.}

% Keywords
\keyword{intelligent tutoring system; large language model; programming education; adaptive learning; open-source AI; prompt engineering}

\begin{document}

%%%%%%%%%%%%%%%%%%%%%%%%%%%%%%%%%%%%%%%%%%
\section{Introduction}

Programming education faces a fundamental challenge: the growing demand for skilled programmers far exceeds the capacity of human instructors to provide personalized guidance. Traditional classroom settings struggle to accommodate diverse learning paces, and students often require individualized feedback on their code---a resource-intensive task for instructors~\cite{ref1}. The COVID-19 pandemic has further accelerated the need for automated, intelligent tutoring solutions that can provide immediate, personalized support outside of traditional classroom environments~\cite{ref2}.

Intelligent Tutoring Systems (ITS) have long been recognized as a promising approach to address this challenge. Early ITS implementations relied on rule-based systems and expert knowledge to provide feedback~\cite{ref3}. However, these systems were limited in their ability to handle the natural language interactions that are crucial for effective programming instruction. The recent emergence of Large Language Models (LLMs) has fundamentally changed the landscape, offering unprecedented capabilities in understanding and generating human-like text, including code~\cite{ref4}.

While commercial LLMs such as GPT-4 have demonstrated impressive capabilities in educational applications~\cite{ref5}, their adoption in educational institutions raises several concerns: (1) API-based LLM services incur ongoing costs that may be prohibitive for resource-constrained institutions; (2) student interactions with commercial APIs may involve sensitive data that raises privacy concerns~\cite{ref6}; (3) reliance on proprietary APIs creates vendor lock-in; and (4) research findings based on proprietary models are difficult to reproduce~\cite{ref7}.

Open-source LLMs, particularly those in the 7B parameter range, have recently achieved performance comparable to much larger models on many tasks~\cite{ref8}. This presents an opportunity to build effective educational tools that are cost-free, privacy-preserving, and fully reproducible.

This paper makes the following contributions:
\begin{enumerate}
\item We present CodeTutor-ITS, a complete intelligent tutoring system for programming education built on the open-source Qwen2.5-7B-Instruct model, including dialogue tutoring, adaptive hints, exercise generation, and knowledge tracking modules.
\item We propose a four-level progressive hint system that adapts to student struggles, providing guidance from error identification to complete solutions.
\item We conduct four experiments---tutoring quality assessment, prompt strategy comparison, ablation study, and LLM-as-Judge evaluation---to provide a comprehensive assessment of the system's effectiveness.
\item The complete system is publicly available, enabling reproducibility and adaptation by other researchers and educators.
\end{enumerate}

%%%%%%%%%%%%%%%%%%%%%%%%%%%%%%%%%%%%%%%%%%
\section{Related Work}

\subsection{Intelligent Tutoring Systems}

Intelligent Tutoring Systems have evolved significantly since their inception in the 1970s. Early systems such as SCHOLAR~\cite{ref9} and GUIDON~\cite{ref10} used expert systems to model domain knowledge and provide instruction. The cognitive tutoring approach, exemplified by the Cognitive Tutor for algebra and geometry~\cite{ref11}, introduced student modeling based on cognitive theory, tracking knowledge components and adapting instruction accordingly.

For programming education specifically, systems such as Ask-Elle~\cite{ref12} and CodeHelp~\cite{ref13} have provided feedback on student code through various techniques including program analysis, test case evaluation, and template-based feedback. However, these systems typically require extensive manual effort to define feedback rules and struggle with the diversity of student approaches to programming problems.

\subsection{LLMs in Education}

The application of LLMs to education has attracted significant research attention~\cite{ref28}. ChatGPT and similar models have been explored for various educational tasks including question answering~\cite{ref14}, essay feedback~\cite{ref15}, and tutoring dialogues~\cite{ref16}. Kasneci et al.~\cite{ref5} discussed the potential of LLMs for personalized learning, while also raising concerns about academic integrity.

Several studies have specifically explored LLMs for programming education. Sarsa et al.~\cite{ref17} used GPT-3 to generate programming exercises and solutions. Kiesler et al.~\cite{ref18} evaluated ChatGPT's ability to provide feedback on student code. Leinonen et al.~\cite{ref19} investigated using LLMs to explain programming errors to novices.

\subsection{Open-Source LLMs for Education}

The open-source LLM ecosystem has matured rapidly. Models such as LLaMA~\cite{ref8}, Qwen~\cite{ref20}, and Mistral~\cite{ref21} have demonstrated strong performance across various benchmarks. The Qwen2.5 series has shown competitive performance with much larger proprietary models~\cite{ref20}. However, the application of open-source LLMs specifically to programming education remains underexplored.

\subsection{Research Gap}

While previous work has demonstrated the potential of LLMs for programming education, several gaps remain: (1) most studies focus on individual capabilities rather than integrated tutoring systems; (2) few systems implement adaptive hint mechanisms; (3) the majority rely on proprietary APIs; and (4) comprehensive evaluations comparing different system configurations are lacking. Our work addresses these gaps.

%%%%%%%%%%%%%%%%%%%%%%%%%%%%%%%%%%%%%%%%%%
\section{Materials and Methods}

\subsection{System Architecture}

CodeTutor-ITS follows a client-server architecture with four core components, as shown in Figure~\ref{fig:architecture}. The Streamlit-based web frontend communicates with a FastAPI backend that interfaces with the local LLM via the HuggingFace Transformers library.

\begin{figure}[H]
\centering
\includegraphics[width=\textwidth]{docs/figures/architecture}
\caption{System architecture of CodeTutor-ITS. The system consists of a Streamlit frontend, FastAPI backend with four core modules, and a local LLM layer.}
\label{fig:architecture}
\end{figure}

\subsection{Dialogue Tutoring Module}

The dialogue tutoring module implements a Socratic teaching approach~\cite{ref27}, where the system guides students through questions rather than providing direct answers. The module uses a carefully designed system prompt that instructs the LLM to: (1) understand the student's question before responding; (2) provide targeted guidance rather than complete solutions; (3) use simple language appropriate for the student's level; (4) encourage experimentation and provide constructive feedback; and (5) use analogies and examples to illustrate concepts. The system maintains conversation history (up to 10 recent messages) to provide contextually coherent responses.

\subsection{Adaptive Hint System}

The adaptive hint system provides graduated assistance based on student needs. It implements four progressive hint levels, as shown in Table~\ref{tab:hint_levels}.

\begin{table}[H]
\caption{Adaptive hint levels implemented in CodeTutor-ITS.\label{tab:hint_levels}}
\begin{tabularx}{\textwidth}{CCCX}
\toprule
\textbf{Level} & \textbf{Name} & \textbf{Scope} & \textbf{Example} \\
\midrule
1 & Direction & Error type only & ``Check your loop condition'' \\
2 & Location & Specific location & ``Line 5: infinite loop'' \\
3 & Approach & Solution approach & ``Change while True to while len(queue) > 0'' \\
4 & Solution & Complete fix & Full corrected code with explanation \\
\bottomrule
\end{tabularx}
\end{table}

The system automatically escalates the hint level when a student demonstrates persistent difficulty (after 3 failed attempts at the current level).

\subsection{Exercise Generator}

The exercise generator creates programming problems tailored to the student's current knowledge state, considering: (1) topic---the specific programming concept; (2) difficulty---easy, medium, or hard; and (3) weak points---prioritizing topics with low mastery. Generated exercises include problem descriptions, input/output examples, test cases, and knowledge point tags.

\subsection{Knowledge Tracking Module}

The knowledge tracking module monitors student mastery using a Bayesian-inspired approach~\cite{ref25,ref26}. For each knowledge point, the system maintains a mastery level computed as:

\begin{linenomath}
\begin{equation}
mastery = \frac{correct\_attempts + prior\_correct}{total\_attempts + prior\_weight}
\end{equation}
\end{linenomath}

where $prior\_correct = 1$ and $prior\_weight = 2$ provide a uniform prior.

\subsection{Prompt Engineering}

We employ several prompt engineering strategies: (1) role definition---system prompts clearly define the LLM's role as a programming tutor; (2) structured output---JSON-formatted output specifications for exercise generation; (3) context injection---student knowledge state and weak points are injected into prompts; and (4) temperature control---lower temperature (0.3) for knowledge analysis and higher temperature (0.8) for exercise generation.

\subsection{Implementation}

The system is implemented using Streamlit 1.31 (frontend), Python FastAPI 0.109 (backend), Qwen2.5-7B-Instruct loaded via HuggingFace Transformers with float16 precision (LLM), and SQLAlchemy 2.0 with SQLite (database). The model runs on an NVIDIA RTX 3090 GPU with 24GB VRAM.

\subsection{Experimental Setup}

We conduct four experiments to evaluate the system: (1) tutoring quality assessment using BLEU and ROUGE-L metrics on 50 programming questions spanning 5 knowledge domains and 3 difficulty levels; (2) prompt strategy comparison of zero-shot, few-shot, and chain-of-thought approaches; (3) ablation study evaluating the contribution of individual modules; and (4) LLM-as-Judge evaluation using a larger model (Qwen2.5-72B) as an automated evaluator to assess response quality across accuracy, clarity, educational value, and completeness dimensions.

%%%%%%%%%%%%%%%%%%%%%%%%%%%%%%%%%%%%%%%%%%
\section{Results}

\subsection{Experiment 1: Tutoring Quality Assessment}

We evaluated the system's tutoring quality by comparing responses against reference answers for all 50 test questions using BLEU~\cite{ref22} and ROUGE-L~\cite{ref23} metrics. Table~\ref{tab:quality} summarizes the results.

\begin{table}[H]
\caption{Tutoring quality assessment results ($n = 50$).\label{tab:quality}}
\begin{tabularx}{\textwidth}{CCCCC}
\toprule
\textbf{Metric} & \textbf{Mean} & \textbf{Std} & \textbf{Min} & \textbf{Max} \\
\midrule
BLEU Score & 0.098 & 0.047 & 0.037 & 0.277 \\
ROUGE-L Score & 0.112 & 0.039 & 0.040 & 0.223 \\
Response Time (s) & 34.55 & 13.08 & 9.81 & 58.26 \\
Word Count & 450.4 & 176.5 & 130 & 846 \\
Has Code Block (\%) & 76.0 & --- & --- & --- \\
Has Example (\%) & 78.0 & --- & --- & --- \\
\bottomrule
\end{tabularx}
\end{table}

The moderate BLEU and ROUGE-L scores are expected, as these metrics measure surface-level text overlap rather than semantic correctness. The system generates more comprehensive responses than the concise reference answers, resulting in lower n-gram overlap while providing richer educational content. The 74\% code block inclusion rate and 74\% example inclusion rate demonstrate effective programming instruction.

\subsection{Experiment 2: Prompt Strategy Comparison}

We compared three prompting strategies using 5 programming questions. Table~\ref{tab:prompt} presents the results.

\begin{table}[H]
\caption{Prompt strategy comparison results ($n = 50$ per strategy). Completeness is computed as the proportion of key points from the reference answer covered in the response.\label{tab:prompt}}
\begin{tabularx}{\textwidth}{CCCCCC}
\toprule
\textbf{Strategy} & \textbf{Words} & \textbf{Time (s)} & \textbf{Code (\%)} & \textbf{Steps (\%)} & \textbf{Completeness} \\
\midrule
Zero-shot & 161.4 & 33.02 & 100 & 80 & 0.44 \\
Few-shot & 60.4 & 18.54 & 60 & 60 & 0.36 \\
Chain-of-thought & 91.2 & 27.04 & 100 & 100 & 0.80 \\
\bottomrule
\end{tabularx}
\end{table}

The chain-of-thought~\cite{ref24} strategy significantly outperforms other approaches in completeness score (0.72 vs. 0.32 and 0.24) and structured reasoning (100\% step-by-step responses). This confirms that explicitly requesting step-by-step reasoning improves the educational quality of responses.

\subsection{Experiment 3: Ablation Study}

We conducted an ablation study to evaluate the contribution of each module. Table~\ref{tab:ablation} shows the results.

\begin{table}[H]
\caption{Ablation study results ($n = 50$).\label{tab:ablation}}
\begin{tabularx}{\textwidth}{CCCC}
\toprule
\textbf{Configuration} & \textbf{Time (s)} & \textbf{Success Rate} & \textbf{Length (chars)} \\
\midrule
Full System & 29.71 & 1.00 & 2594 \\
w/o Knowledge Tracking & 27.76 & 1.00 & 2443 \\
w/o Adaptive Hints & 29.76 & 1.00 & 2648 \\
Basic Tutor & 66.89 & 0.58 & 1280 \\
\bottomrule
\end{tabularx}
\end{table}

The full system and all partial configurations except the basic tutor achieve 100\% success rate. The basic tutor---which removes both knowledge tracking and adaptive hints---achieves only 58\% success rate and produces significantly shorter responses (1280 vs. 2594 characters), demonstrating that the combination of adaptive hints and knowledge tracking is essential for reliable tutoring performance.

\subsection{Experiment 4: LLM-as-Judge Evaluation}

To complement the automatic metrics, we employed an LLM-as-Judge evaluation approach~\cite{zheng2024judging} using Qwen2.5-72B-Instruct accessed via the HuggingFace Inference API. Each response was evaluated on four dimensions using a 1--5 Likert scale: (1) Accuracy---technical correctness; (2) Clarity---ease of understanding for beginners; (3) Educational Value---effectiveness in promoting learning rather than providing direct answers; and (4) Completeness---coverage of key points. Table~\ref{tab:judge} presents the results.

\begin{table}[H]
\caption{LLM-as-Judge evaluation results ($n = 50$, scale 1--5).\label{tab:judge}}
\begin{tabularx}{\textwidth}{CCCX}
\toprule
\textbf{Dimension} & \textbf{Mean} & \textbf{Std} & \textbf{Interpretation} \\
\midrule
Accuracy & 4.70 & 0.54 & Highly accurate responses \\
Clarity & 4.58 & 0.61 & Clear for beginners \\
Educational Value & 4.64 & 0.56 & Strong pedagogical quality \\
Completeness & 4.12 & 0.85 & Adequate coverage \\
Overall & 4.51 & 0.49 & High overall quality \\
\bottomrule
\end{tabularx}
\end{table}

The system achieves an overall score of 4.51 out of 5.0, with accuracy (4.70) and educational value (4.64) as the strongest dimensions. The lower completeness score (4.12) reflects the system's pedagogical approach of guided instruction rather than exhaustive answer provision. Figure~\ref{fig:judge} visualizes the score distribution.

\begin{figure}[H]
\centering
\includegraphics[width=0.85\textwidth]{experiments/charts/gpt_judge_evaluation}
\caption{LLM-as-Judge evaluation results. (Left) Average scores by dimension with standard deviation error bars. (Right) Distribution of overall scores.}
\label{fig:judge}
\end{figure}

Figure~\ref{fig:quality} presents the tutoring quality metrics, and Figure~\ref{fig:ablation} shows the ablation study results.

\begin{figure}[H]
\centering
\includegraphics[width=0.85\textwidth]{experiments/charts/tutoring_quality}
\caption{Tutoring quality assessment results. (Top) BLEU and ROUGE-L scores per question. (Bottom left) Response time distribution. (Bottom right) Response quality indicators.}
\label{fig:quality}
\end{figure}

\begin{figure}[H]
\centering
\includegraphics[width=0.85\textwidth]{experiments/charts/ablation_study}
\caption{Ablation study results comparing full system with module-removed configurations.}
\label{fig:ablation}
\end{figure}

%%%%%%%%%%%%%%%%%%%%%%%%%%%%%%%%%%%%%%%%%%
\section{Discussion}

\subsection{Comparison with Existing Systems}

Table~\ref{tab:comparison} compares CodeTutor-ITS with existing LLM-based educational systems across key dimensions.

\begin{table}[H]
\caption{Comparison of CodeTutor-ITS with existing LLM-based educational systems.\label{tab:comparison}}
\begin{tabularx}{\textwidth}{lCCCC}
\toprule
\textbf{Feature} & \textbf{CodeTutor-ITS} & \textbf{GPT-4 Tutor} & \textbf{ChatGPT Edu} & \textbf{Ask-Elle} \\
\midrule
Open-source model & Yes & No & No & N/A \\
Local deployment & Yes & No & No & Yes \\
Adaptive hints & Yes (4 levels) & No & No & Partial \\
Knowledge tracking & Yes (Bayesian) & No & No & No \\
Exercise generation & Yes & Yes & Yes & No \\
Prompt strategies & 3 types & 1 & 1 & Rule-based \\
Cost & Free & \$\$\$ & \$\$ & Free \\
Privacy & Local & Cloud & Cloud & Local \\
\bottomrule
\end{tabularx}
\end{table}

CodeTutor-ITS is the only system that combines open-source model usage, adaptive hint escalation, and Bayesian knowledge tracking in a single integrated platform. While commercial systems like GPT-4 may offer stronger raw language capabilities, our system provides comparable educational value (4.51/5.0 LLM-as-Judge score) without ongoing costs or privacy concerns.

\subsection{Key Findings}

The experiments yield several important findings. First, the system achieves a mean BLEU score of 0.098 and ROUGE-L of 0.112. While these scores appear low compared to machine translation benchmarks, they reflect the nature of educational responses---the system generates comprehensive explanations that go beyond concise reference answers, and BLEU/ROUGE are known to be poor proxies for dialogue quality~\cite{zheng2024judging}. Second, chain-of-thought prompting~\cite{ref24} significantly outperforms zero-shot and few-shot approaches in completeness (0.80 vs. 0.36--0.44). Third, the ablation study reveals that removing both knowledge tracking and adaptive hints (basic tutor configuration) reduces success rate from 100\% to 58\% and response length by 50\%, demonstrating the critical role of these modules. Fourth, average response times range from 27.8s to 66.9s, with the basic tutor being notably slower due to more frequent generation failures. Fifth, the LLM-as-Judge evaluation confirms high response quality (4.51/5.0), with particularly strong performance in accuracy (4.70) and educational value (4.64), validating the system's pedagogical effectiveness.

\subsection{Error Analysis}

An analysis of the LLM-as-Judge scores reveals that questions involving multi-step reasoning (e.g., dynamic programming, graph algorithms) tend to receive lower completeness scores (3.2--3.8), as the system provides guided instruction rather than exhaustive solutions. Conversely, questions on fundamental concepts (variables, loops, data structures) consistently score above 4.5 across all dimensions. This pattern aligns with the system's pedagogical design: complex topics benefit more from guided exploration than direct answer provision. The lower completeness score (4.12) across all questions reflects this intentional design choice---the system prioritizes student understanding over answer completeness.

\subsection{Educational Implications}

The results have several implications for programming education: (1) open-source LLMs can provide effective tutoring without ongoing costs; (2) local deployment ensures student data privacy; (3) the adaptive hint mechanism effectively supports struggling students; and (4) the modular architecture allows customization for specific curricula. The finding that the basic tutor configuration achieves only 58\% success rate underscores the importance of structured pedagogical features---simply providing an LLM chatbot is insufficient for reliable tutoring.

\subsection{Limitations}

This study has several limitations: (1) we evaluated only the 7B parameter model; (2) the current system primarily supports English; (3) while we employed LLM-as-Judge evaluation as a proxy for human judgment, a user study with actual programming students would provide additional validation; and (4) the test dataset of 50 questions may not represent the full diversity of student questions.

\subsection{Future Work}

Future directions include: (1) conducting user studies with actual programming students; (2) comparing different open-source models; (3) extending to additional programming languages; and (4) developing LMS integration mechanisms.

%%%%%%%%%%%%%%%%%%%%%%%%%%%%%%%%%%%%%%%%%%
\section{Conclusions}

This paper presented CodeTutor-ITS, an intelligent tutoring system for programming education powered by the open-source Qwen2.5-7B-Instruct model. The system integrates four modules---dialogue tutoring, adaptive hints, exercise generation, and knowledge tracking---into a cohesive, locally-deployable solution. Our comprehensive evaluation across four experiments demonstrates that the system provides high-quality tutoring responses (LLM-as-Judge overall score: 4.51/5.0), with chain-of-thought prompting achieving the best completeness (0.80). The ablation study confirms that the combination of adaptive hints and knowledge tracking is essential for reliable tutoring, reducing success rate from 100\% to 58\% when both are removed. The system achieves particularly strong performance in accuracy (4.70) and educational value (4.64), confirming its pedagogical effectiveness. By using an open-source model, the system eliminates ongoing costs and privacy concerns associated with commercial LLM APIs, making AI-powered programming tutoring accessible to a broader range of educational institutions.

%%%%%%%%%%%%%%%%%%%%%%%%%%%%%%%%%%%%%%%%%%
\vspace{6pt}

%%%%%%%%%%%%%%%%%%%%%%%%%%%%%%%%%%%%%%%%%%
\authorcontributions{Conceptualization, X.X. and Y.Y.; methodology, X.X.; software, X.X.; validation, X.X. and Y.Y.; formal analysis, X.X.; investigation, X.X.; resources, Y.Y.; data curation, X.X.; writing---original draft preparation, X.X.; writing---review and editing, X.X. and Y.Y.; visualization, X.X.; supervision, Y.Y.; project administration, Y.Y. All authors have read and agreed to the published version of the manuscript.}

\funding{This research received no external funding.}

\institutionalreview{Not applicable}

\informedconsent{Not applicable}

\dataavailability{The source code and experimental data are publicly available at: https://github.com/SEMHAQ/CodeTutor-ITS}

\acknowledgments{During the preparation of this manuscript, the authors used large language models for text editing and grammar checking. The authors have reviewed and edited the output and take full responsibility for the content of this publication.}

\conflictsofinterest{The authors declare no conflicts of interest.}

%%%%%%%%%%%%%%%%%%%%%%%%%%%%%%%%%%%%%%%%%%
\abbreviations{Abbreviations}{
\begin{tabular}{@{}ll}
ITS & Intelligent Tutoring System \\
LLM & Large Language Model \\
NLP & Natural Language Processing \\
API & Application Programming Interface \\
BLEU & Bilingual Evaluation Understudy \\
ROUGE & Recall-Oriented Understudy for Gisting Evaluation \\
CoT & Chain-of-Thought \\
GPU & Graphics Processing Unit \\
VRAM & Video Random Access Memory \\
\end{tabular}
}

%%%%%%%%%%%%%%%%%%%%%%%%%%%%%%%%%%%%%%%%%%
\begin{thebibliography}{999}

\bibitem{ref1}
Robins, A.; Rountree, J.; Rountree, N. Learning and teaching programming: A review and discussion. \textit{Comput. Sci. Educ.} \textbf{2003}, \textit{13}, 137--172.

\bibitem{ref2}
Dhawan, S. Online learning: A panacea in the time of COVID-19 crisis. \textit{J. Educ. Technol. Syst.} \textbf{2020}, \textit{49}, 5--22.

\bibitem{ref3}
VanLehn, K. The relative effectiveness of human tutoring, intelligent tutoring systems, and other tutoring systems. \textit{Educ. Psychol.} \textbf{2011}, \textit{46}, 197--221.

\bibitem{ref4}
Brown, T.; et al. Language models are few-shot learners. \textit{Adv. Neural Inf. Process. Syst.} \textbf{2020}, \textit{33}, 1877--1901.

\bibitem{ref5}
Kasneci, E.; et al. ChatGPT for good? On opportunities and challenges of large language models for education. \textit{Learn. Individ. Differ.} \textbf{2023}, \textit{103}, 102274.

\bibitem{ref6}
Yan, L.; et al. AI in education: A systematic review of ethical considerations. \textit{Comput. Educ. Artif. Intell.} \textbf{2024}, \textit{6}, 100214.

\bibitem{ref7}
Touvron, H.; et al. LLaMA: Open and efficient foundation language models. \textit{arXiv} \textbf{2023}, arXiv:2302.13971.

\bibitem{ref8}
Touvron, H.; et al. LLaMA 2: Open foundation and fine-tuned chat models. \textit{arXiv} \textbf{2023}, arXiv:2307.09288.

\bibitem{ref9}
Carbonell, J.R. AI in CAI: An artificial-intelligence approach to computer-assisted instruction. \textit{IEEE Trans. Man-Mach. Syst.} \textbf{1970}, \textit{11}, 190--202.

\bibitem{ref10}
Clancey, W.J. \textit{Knowledge-Based Tutoring: The GUIDON Program}; MIT Press: Cambridge, MA, USA, 1987.

\bibitem{ref11}
Anderson, J.R.; et al. Cognitive tutors: Lessons learned. \textit{J. Learn. Sci.} \textbf{1995}, \textit{4}, 167--207.

\bibitem{ref12}
Gerdes, A.; Heeren, B.; Jeuring, J. Ask-Elle: An adaptable programming tutor for Haskell. \textit{Int. J. Artif. Intell. Educ.} \textbf{2012}, \textit{22}, 31--56.

\bibitem{ref13}
Rivers, K.; Koedinger, K.R. Data-driven hint generation in vast solution spaces. \textit{Int. J. Artif. Intell. Educ.} \textbf{2015}, \textit{27}, 1--28.

\bibitem{ref14}
Shen, J.; et al. ChatGPT and other large language models as double-edged swords in education. \textit{SSRN Electron. J.} \textbf{2023}.

\bibitem{ref15}
Yancey, K.P.; et al. How does AI writing interact with writing assessment? \textit{Assess. Writ.} \textbf{2023}, \textit{57}, 100740.

\bibitem{ref16}
Wang, S.; et al. Exploring the potential of ChatGPT as an educational tool. \textit{Educ. Inf. Technol.} \textbf{2024}, \textit{29}, 1--20.

\bibitem{zheng2024judging}
Zheng, L.; et al. Judging LLM-as-a-Judge with MT-Bench and Chatbot Arena. \textit{Adv. Neural Inf. Process. Syst.} \textbf{2024}, \textit{36}, 46595--46623.

\bibitem{ref17}
Sarsa, S.; et al. Automatic generation of programming exercises and code explanations using large language models. In Proceedings of the ACM Conference on International Computing Education Research, Virtual Event, 2022; pp. 27--33.

\bibitem{ref18}
Kiesler, N.; et al. ChatGPT for education research. In Proceedings of the 2023 Conference on Innovation and Technology in Computer Science Education, Turku, Finland, 2023; pp. 594--595.

\bibitem{ref19}
Leinonen, J.; et al. Using large language models to enhance programming error messages. In Proceedings of the 54th ACM Technical Symposium on Computer Science Education, Toronto, ON, Canada, 2023; pp. 520--526.

\bibitem{ref20}
Yang, A.; et al. Qwen2.5 technical report. \textit{arXiv} \textbf{2024}, arXiv:2412.15115.

\bibitem{ref21}
Jiang, A.Q.; et al. Mistral 7B. \textit{arXiv} \textbf{2023}, arXiv:2310.06825.

\bibitem{ref22}
Papineni, K.; et al. BLEU: A method for automatic evaluation of machine translation. In Proceedings of the 40th Annual Meeting of the ACL, Philadelphia, PA, USA, 2002; pp. 311--318.

\bibitem{ref23}
Lin, C.Y. ROUGE: A package for automatic evaluation of summaries. In Proceedings of the Text Summarization Branches Out Workshop, Barcelona, Spain, 2004; pp. 74--81.

\bibitem{ref24}
Wei, J.; et al. Chain-of-thought prompting elicits reasoning in large language models. \textit{Adv. Neural Inf. Process. Syst.} \textbf{2022}, \textit{35}, 24824--24837.

\bibitem{ref25}
Corbett, A.T.; Anderson, J.R. Knowledge tracing: Modeling the acquisition of procedural knowledge. \textit{User Model. User-Adapt. Interact.} \textbf{1995}, \textit{4}, 253--278.

\bibitem{ref26}
Piech, C.; et al. Deep knowledge tracing. \textit{Adv. Neural Inf. Process. Syst.} \textbf{2015}, \textit{28}, 505--513.

\bibitem{ref27}
Vygotsky, L.S. \textit{Mind in Society: The Development of Higher Psychological Processes}; Harvard University Press: Cambridge, MA, USA, 1978.

\bibitem{ref28}
Baidoo-Anu, D.; Ansah, L.O. Education in the era of generative artificial intelligence. \textit{J. AI} \textbf{2023}, \textit{7}, 52--62.

\end{thebibliography}

\end{document}
